\documentclass[10pt,aspectratio=169]{beamer}
% Preámbulo modular para presentaciones Beamer (Modelación Estadística)
% Diseño y paleta institucional del Tecnológico de Monterrey

\documentclass[aspectratio=169,xcolor={dvipsnames,table}]{beamer}

\usepackage[utf8]{inputenc}
\usepackage[spanish,mexico]{babel}
\usepackage{amsmath,amssymb,amsfonts,mathtools}
\usepackage{graphicx}
\usepackage{listings}
\usepackage{booktabs}
\usepackage{tikz}
\usetikzlibrary{matrix,arrows,decorations.pathmorphing,shapes.geometric,positioning}

% Paleta Institucional Tecnológico de Monterrey
\definecolor{TecAzul}{HTML}{0039A6}       % Azul TEC: color primario institucional
\definecolor{TecAzulOscuro}{HTML}{1A2E51} % Azul marino oscuro
\definecolor{TecRojo}{HTML}{EC2661}       % Rojo de acento (paleta miTec)
\definecolor{TecGris}{HTML}{646464}       % Gris neutro para comentarios y subtítulos
\definecolor{TecFondoBloque}{HTML}{F4F6F9}% Fondo suave para bloques

% Configuración de Tema Beamer
\usetheme{Boadilla}
\usecolortheme[named=TecAzulOscuro]{structure}
\setbeamercolor{alerted text}{fg=TecRojo}
\setbeamercolor{palette primary}{bg=TecAzulOscuro,fg=white}
\setbeamercolor{palette secondary}{bg=TecAzul,fg=white}
\setbeamercolor{palette tertiary}{bg=TecAzulOscuro!80!black,fg=white}
\setbeamercolor{title}{fg=TecAzulOscuro,bg=TecFondoBloque}
\setbeamercolor{frametitle}{fg=TecAzulOscuro,bg=TecFondoBloque}

% Bloques personalizados elegantes
\setbeamercolor{block title}{bg=TecAzulOscuro,fg=white}
\setbeamercolor{block body}{bg=TecFondoBloque,fg=black}
\setbeamercolor{block title alerted}{bg=TecRojo,fg=white}
\setbeamercolor{block body alerted}{bg=TecRojo!8,fg=black}
\setbeamercolor{block title example}{bg=TecAzul,fg=white}
\setbeamercolor{block body example}{bg=TecAzul!8,fg=black}

% Redondear bordes y añadir sombra sutil
\setbeamertemplate{blocks}[rounded][shadow=true]
\setbeamertemplate{navigation symbols}{}

% Pie de página limpio con número de diapositiva
\setbeamertemplate{footline}{
  \leavevmode%
  \hbox{%
  \begin{beamercolorbox}[wd=.7\paperwidth,ht=2.25ex,dp=1ex,left]{author in head/foot}%
    \hspace*{2ex}\insertshortauthor~~(\insertshortinstitute) \hfill \insertshorttitle
  \end{beamercolorbox}%
  \begin{beamercolorbox}[wd=.3\paperwidth,ht=2.25ex,dp=1ex,right]{date in head/foot}%
    \insertshortdate{}\hspace*{2em}
    \insertframenumber{} / \inserttotalframenumber\hspace*{2ex}
  \end{beamercolorbox}}%
  \vskip0pt%
}

% Configuración de Listings de Python para visualización pedagógica de código
\lstset{
    language=Python,
    basicstyle=\ttfamily\footnotesize,
    keywordstyle=\color{TecAzulOscuro}\bfseries,
    stringstyle=\color{TecRojo},
    commentstyle=\color{TecGris}\itshape,
    showstringspaces=false,
    breaklines=true,
    frame=lines,
    rulecolor=\color{TecAzulOscuro!30},
    backgroundcolor=\color{TecFondoBloque},
    aboveskip=0.5em,
    belowskip=0.5em,
    literate={á}{{\'a}}1 {é}{{\'e}}1 {í}{{\'i}}1 {ó}{{\'o}}1 {ú}{{\'u}}1
             {Á}{{\'A}}1 {É}{{\'E}}1 {Í}{{\'I}}1 {Ó}{{\'O}}1 {Ú}{{\'U}}1
             {ñ}{{\~n}}1 {Ñ}{{\~N}}1 {¿}{{?`}}1 {¡}{{!`}}1
}

% Metadatos institucionales por defecto
\title[Modelación Estadística]{Modelación Estadística para Ciencia de Datos e Ingeniería}
\author[J. Castillo Colmenares]{Juliho David Castillo Colmenares}
\institute[Tec de Monterrey]{Tecnológico de Monterrey \\ Escuela de Ingeniería y Ciencias}
\date{\today}
% Comandos y macros estadísticas compartidas para las presentaciones Beamer (Metropolis)

% Conjuntos numéricos básicos
\newcommand{\R}{\mathbb{R}}
\newcommand{\N}{\mathbb{N}}
\newcommand{\Q}{\mathbb{Q}}
\newcommand{\Z}{\mathbb{Z}}
\newcommand{\set}[1]{\left\{ #1 \right\}}
\newcommand{\abs}[1]{\left|#1\right|}
\newcommand{\norm}[1]{\left\| #1 \right\|}

% Comandos estadísticos y probabilísticos del curso
\newcommand{\Var}{\operatorname{Var}}
\newcommand{\cov}{\operatorname{Cov}}
\newcommand{\corr}{\rho}
\newcommand{\comb}[2]{\begin{pmatrix} #1 \\ #2 \end{pmatrix}}
\newcommand{\s}{\sigma}
\newcommand{\particion}{\mathfrak{P}}
\newcommand{\card}[1]{\#\left( #1 \right)}
\newcommand{\seq}[3]{\set{#1}_{#2}^{#3}}
\newcommand{\E}{\mathbb{E}}
\newcommand{\Prob}{\mathbb{P}}

% Entornos pedagógicos y cajas en español académico natural
\newenvironment{discusion}{%
  \begin{alertblock}{Pregunta de Discusión}%
}{%
  \end{alertblock}%
}

\newenvironment{reflexion}{%
  \begin{alertblock}{Reflexión para el Aula}%
}{%
  \end{alertblock}%
}

\newenvironment{paradoja}{%
  \begin{alertblock}{Paradoja de Exploración}%
}{%
  \end{alertblock}%
}

\newenvironment{motivacion}{%
  \begin{exampleblock}{Motivación: Ciencia de Datos en la Práctica}%
}{%
  \end{exampleblock}%
}

\newenvironment{retoclase}{%
  \begin{block}{Reto Rápido en Equipo}%
}{%
  \end{block}%
}

\title[Pruebas de proporciones]{Sección 07.07: Pruebas relacionadas con proporciones}\subtitle{Una y dos poblaciones binarias}
\author[J. Castillo Colmenares]{Juliho Castillo Colmenares}\institute{Tecnológico de Monterrey}\date{}
\begin{document}
\begin{frame}[plain]\titlepage\end{frame}
\begin{frame}{Hoja de ruta}\small\begin{enumerate}\item Una proporción.\item Diferencia de dos proporciones.\item Condiciones y decisiones.\end{enumerate}\end{frame}
\begin{frame}{Fuente y alcance}\small La sección contrasta proporciones de éxito/fracaso en variables categóricas o binarias.\begin{block}{Pregunta guía}¿La proporción observada es compatible con un valor de referencia?\end{block}\end{frame}
\begin{frame}{Modelo binomial}\small Sea $X$ el número de éxitos en $n$ ensayos, con probabilidad $p$. La proporción muestral es $\hat p=X/n$.\end{frame}
\begin{frame}{Hipótesis para una proporción}\small\[H_0:p=p_0\qquad\text{frente a una alternativa sobre }p.\]\end{frame}
\begin{frame}{Estadística para una proporción}\small\[Z=\frac{\hat p-p_0}{\sqrt{p_0(1-p_0)/n}}.\]\end{frame}
\begin{frame}{Aproximación normal}\small La aproximación requiere $np_0\geq5$ y $n(1-p_0)\geq5$. Si no se cumplen, se debe considerar un método exacto.\end{frame}
\begin{frame}{Prueba de dos proporciones}\small Para dos poblaciones independientes, la hipótesis de igualdad es $H_0:p_1-p_2=0$. Se observan $\hat p_1$ y $\hat p_2$.\end{frame}
\begin{frame}{Proporción combinada}\small Bajo $H_0:p_1=p_2=p$, se estima\[\bar p=\frac{X_1+X_2}{n_1+n_2}.\]\end{frame}
\begin{frame}{Estadística de diferencia}\small\[Z=\frac{\hat p_1-\hat p_2}{\sqrt{\bar p(1-\bar p)(1/n_1+1/n_2)}}.\]\end{frame}
\begin{frame}{Procedimiento I}\small\begin{enumerate}\item Definir éxito y fracaso.\item Formular hipótesis.\item Verificar condiciones de aproximación.\end{enumerate}\end{frame}
\begin{frame}{Procedimiento II}\small\begin{enumerate}\setcounter{enumi}{3}\item Calcular proporciones y estadística.\item Obtener valor crítico o $p$-value.\item Comunicar la decisión.\end{enumerate}\end{frame}
\begin{frame}{Dirección de la alternativa}\small Las colas dependen de si se busca una proporción mayor, menor o distinta. El signo de $\hat p-p_0$ orienta la interpretación.\end{frame}
\begin{frame}{Aplicación: una proporción}\small Una muestra contiene $X$ éxitos de $n$ observaciones. Se contrasta $H_0:p=p_0$ al nivel $\alpha$.\end{frame}
\begin{frame}{Aplicación: cálculo}\small Calcule $\hat p=X/n$ y sustituya en $Z$. El denominador se evalúa bajo $p_0$, no bajo $\hat p$.\end{frame}
\begin{frame}{Aplicación: decisión}\small Compare $Z$ con la región crítica de la alternativa o use el valor $p$.\end{frame}
\begin{frame}{Aplicación: dos poblaciones}\small Dos muestras independientes aportan $X_1,n_1$ y $X_2,n_2$. Se contrasta igualdad mediante la proporción combinada.\end{frame}
\begin{frame}{Aplicación: interpretación}\small Una diferencia estadísticamente detectable entre $p_1$ y $p_2$ debe reportarse junto con sus proporciones, tamaños muestrales y dirección.\end{frame}
\begin{frame}{Errores frecuentes}\small\begin{itemize}\item Usar $\hat p$ en el error nulo.\item Ignorar condiciones de aproximación.\item Confundir diferencia estadística con importancia práctica.\end{itemize}\end{frame}
\begin{frame}{Plantilla de reporte}\small\begin{block}{Reporte} ``$\hat p=\cdots$ (o $\hat p_1-\hat p_2=\cdots$), $Z=\cdots$, $p=\cdots$. Al nivel $\alpha=\cdots$, [rechazamos/no rechazamos] $H_0$.''\end{block}\end{frame}
\begin{frame}{Síntesis}\small\begin{itemize}\item Una proporción usa el valor nulo $p_0$.\item Dos proporciones usan una estimación combinada bajo igualdad.\item Las condiciones sostienen la aproximación normal.\end{itemize}\end{frame}
\begin{frame}{Conexión con el capítulo}\small Estas pruebas preparan la comparación de varias proporciones y las pruebas de homogeneidad.\end{frame}
\end{document}
